We now consider the two possible cases giving equation(3.5).
\begin{enumerate}
\item Case1 explaining equation(3.5): After rescaling $C$ and $D$, we have that $B''=C''=D''$. Equation(3.2) implies that for some linear form $u,v$ we have that
\[C=bv+B'\quad\text{and}\quad D=bu+B'.\]
We now expand the resultant again:
\begin{align*}B'^2+b(v+u)B'+b^2vu&=(bv+B')\cdot(bu+B')=CD\\&=\operatorname{Res}_{x_1}(A,B)=B'^2-b^2A'.\end{align*}
Hence,
\begin{equation}(v+u)B'+bvu=-bA'.\tag{3.7}\end{equation}
Thus, either $b$ divides $B'$ in which case we get that $b$ divides $B$ and we are done as Theorem3.1(product-case) holds, or $b$ divides $u+v$. That is,
\begin{equation}u+v=\varepsilon b\tag{3.8}\end{equation}
for some constant $\varepsilon\in\mathbb C$. Plugging this back into equation(3.7), we get
\[\varepsilon bB'+bvu=-bA'.\]
In other words, $\varepsilon B'+vu=-A'$.
Consider the linear combination $Q=A+\varepsilon B$. We get that
\begin{align}Q=A+\varepsilon B&=(x_1^2-A')+\varepsilon(x_1b-B')\notag\\&=x_1^2+\varepsilon x_1b+vu\notag\\&=x_1^2+x_1(u+v)+uv\notag\\&=(x_1+u)(x_1+v),\tag{3.9}\end{align}
where the equality in the third line follows from equation(3.8). Thus, equation(3.9) shows that some linear combination of $A$ and $B$ is reducible which implies that Theorem3.1(product-case) holds.
\item Case2 explaining equation(3.5): $B''=u\cdot v$ and we have that, without loss of generality, $C''=u^2$ and $D''=v^2$ (where $u,v$ are linear forms). Consider equation(3.6). We have that $v$ divides $2a_2B''-a_3D''$. It follows that $v$ is also a factor of $a_4C''$. Thus, either $u$ is a multiple of $v$ and we are back in the case where $C''$ and $D''$ are multiples of each other, or $a_4$ is a multiple of $v$. In this case, we get from equation(3.2) that for some constant $\delta'$,
\[D=\delta b^2+a_4b+D''=\delta b^2+\delta'vb+v^2.\]
Thus, $D$ is a homogeneous polynomial in two linear forms. Hence, $D$ is reducible, in contradiction.
\end{enumerate}
This concludes the proof of Theorem3.1 for the case $b\not\equiv0$.
\paragraph{Case $b\equiv0$}
To ease notation, let use denote $x=x_1$. We have that $A=x^2-A'$ and that $x$ does not appear in $A',B$. Let $y$ be some variable such that $B=y^2-B'$, and $B'$ does not involve $y$ (we can always assume this is the case without loss of generality). As before, we can subtract a multiple of $B$ from $A$ so that the term $y^2$ does not appear in $A$. If $A$ still involves $y$, then we are back in the previous case (treating $y$ as the variable according to which we take the resultant). Thus, the only case left to study is when there are two variables $x$ and $y$ such that
\[A=x^2-A'\quad\text{and}\quad B=y^2-B',\]
where neither $A'$ nor $B'$ involve either $x$ or $y$. To ease notation, denote the rest of the variables as $\mathbf z$. Thus, $A'=A'(\mathbf z)$ and $B'=B'(\mathbf z)$. It is immediate that for any assignment to $\mathbf z$ there is an assignment to $x,y$ that yields a common zero of $A,B$.
By subtracting linear combinations of $A$ and $B$ from the $Q_i$s, we can assume that for every $i\in[4]$
\[Q_i=\alpha_i xy+a_i(\mathbf z)x+b_i(\mathbf z)y+Q'_i(\mathbf z).\]
We next show that, under the assumptions in the theorem's statement, it must be the case that either $A'$ or $B'$ is a perfect square or that $A'\sim B'$. In either situation, we have that Theorem3.1(product-case) holds. We first show that, if $A'$ and $B'$ are linearly independent, then this implies that at least one of $A',B'$ is a perfect square.
Let $Z(A,B)$ be the set of common zeros of $A$ and $B$, and denote by $\pi_{\mathbf z}:Z(A,B)\to\mathbb C^{n-2}$, the projection on the $\mathbf z$ coordinates. Note that $\pi_{\mathbf z}$ is surjective; as for any assignment to $\mathbf z$, there is an assignment to $x,y$ that yields a common zero of $A,B$.
\begin{ntclaim}[Author Claim3.4]
Let $Z(A,B)=\bigcup_{i=1}^kX_k$ be the decomposition of $Z(A,B)$ to irreducible components. Then there exists $i\in[k]$ such that $\pi_{\mathbf z}(X_i)$ is dense in $\mathbb C^{n-2}$.
\end{ntclaim}
\begin{proof}
$\bigcup_{i=1}^k\pi_{\mathbf z}(X_i)=\pi_{\mathbf z}(Z(A,B))=\mathbb C^{n-2}$, as $\pi_{\mathbf z}$ is a surjection, it holds that $\bigcup_{i=1}^k\overline{\pi_{\mathbf z}(X_i)}=\mathbb C^{n-2}$. We also know that $\mathbb C^{n-2}$ is irreducible, and thus, there is $i\in[k]$ such that $\overline{\pi_{\mathbf z}(X_i)}=\mathbb C^{n-2}$, which implies that $\pi_{\mathbf z}(X_i)$ is dense.
\end{proof}
Assume, without loss of generality, that $\pi_{\mathbf z}(X_1)$ is dense. We know that $X_1\subseteq Z(\prod_{i=1}^4Q_i)$, so we can assume, without loss of generality, that $X_1\subseteq Z(Q_1)$. Observe that this implies that $Q_1$ must depend on at least one of $x,y$. Indeed, if $Q_1$ depends on neither, then it is a polynomial in $\mathbf z$, and hence, its set of zeros cannot be dense.
Every point $\boldsymbol\xi\in X_1$ is of the form $\boldsymbol\xi=(\delta_1\sqrt{A'(\boldsymbol\beta)},\delta_2\sqrt{B'(\boldsymbol\beta)},\boldsymbol\beta)$, for some $\boldsymbol\beta\in\mathbb C^{n-2}$, $\delta_1,\delta_2\in\{\pm1\}$ ($\delta_1,\delta_2$ may be functions of $\boldsymbol\beta$). Thus, $Q_1(\boldsymbol\xi)=Q_1(\delta_1\sqrt{A'(\boldsymbol\beta)},\delta_2\sqrt{B'(\boldsymbol\beta)},\boldsymbol\beta)=0$, and we obtain that
\begin{equation}\alpha_1\delta_1\delta_2\sqrt{A'(\boldsymbol\beta')}\sqrt{B'(\boldsymbol\beta')}+a_1(\boldsymbol\beta')\delta_1\sqrt{A'(\boldsymbol\beta')}+b_1(\boldsymbol\beta')\delta_2\sqrt{B'(\boldsymbol\beta')}+Q'_1(\boldsymbol\beta')=0.\tag{3.10}\end{equation}
As we assumed that $Q_1$ depends on at least one of $x,y$, let us assume without loss of generality that either $\alpha_1$ or $a_1$ are nonzero. The next argument is similar to the proof that $\sqrt2$ is irrational. Note that we use the fact that $\delta_1^2=\delta_2^2=1$.
\begin{align*}(3.10)\ \Longrightarrow\quad&B'(\boldsymbol\beta')\left(\alpha_1\delta_1\sqrt{A'(\boldsymbol\beta')}+b_1(\boldsymbol\beta')\right)^2=\left(Q'_1(\boldsymbol\beta')+a_1(\boldsymbol\beta')\delta_1\sqrt{A'(\boldsymbol\beta')}\right)^2\\\Longrightarrow\quad&B'(\boldsymbol\beta')\left(\alpha_1^2A'(\boldsymbol\beta')+2\delta_1\alpha_1b_1(\boldsymbol\beta')\sqrt{A'(\boldsymbol\beta')}+b_1(\boldsymbol\beta')^2\right)\\&\quad=Q'_1(\boldsymbol\beta')^2+2\delta_1a_1(\boldsymbol\beta')Q'_1(\boldsymbol\beta')\sqrt{A'(\boldsymbol\beta')}+a_1(\boldsymbol\beta')^2A'(\boldsymbol\beta').\end{align*}
\begin{align}\Longrightarrow\quad&\delta_1\sqrt{A'(\boldsymbol\beta')}\left(2\alpha_1b_1(\boldsymbol\beta')B'(\boldsymbol\beta')-2a_1(\boldsymbol\beta')Q'_1(\boldsymbol\beta')\right)\notag\\&\quad=Q'_1(\boldsymbol\beta')^2+a_1(\boldsymbol\beta')^2A'(\boldsymbol\beta')-B'(\boldsymbol\beta')\left(\alpha_1^2A'(\boldsymbol\beta')+b_1(\boldsymbol\beta')^2\right)\tag{3.11}\\\Longrightarrow\quad&A'(\boldsymbol\beta')\left(2\alpha_1b_1(\boldsymbol\beta')B'(\boldsymbol\beta')-2a_1(\boldsymbol\beta')Q'_1(\boldsymbol\beta')\right)^2\notag\\&\quad=\left(Q'_1(\boldsymbol\beta')^2+a_1(\boldsymbol\beta')^2A'(\boldsymbol\beta')-B'(\boldsymbol\beta')\left(\alpha_1^2A'(\boldsymbol\beta')+b_1(\boldsymbol\beta')^2\right)\right)^2.\tag{3.12}\end{align}
This equality holds for every $\boldsymbol\beta\in\pi_{\mathbf z}(X_1)$, which is a dense set, and hence holds as a polynomial identity. Thus, either $A'(\mathbf z)$ is a square, in which case we are done or it must be the case that the following identities hold
\begin{equation}Q'_1(\mathbf z)^2+a_1(\mathbf z)^2A'(\mathbf z)-B'(\mathbf z)\left(\alpha_1^2A'(\mathbf z)+b_1(\mathbf z)^2\right)=0\tag{3.13}\end{equation}
and
\begin{equation}\alpha_1b_1(\mathbf z)B'(\mathbf z)-a_1(\mathbf z)Q'_1(\mathbf z)=0.\tag{3.14}\end{equation}
By symmetry, if $B'(\mathbf z)$ is not a square (as otherwise we are done), we get that
\begin{equation}\alpha_1a_1(\mathbf z)A'(\mathbf z)-b_1(\mathbf z)Q'_1(\mathbf z)=0.\tag{3.15}\end{equation}
If $\alpha_1=0$, then we get from(3.14) that $Q'_1\equiv0$. Hence, by(3.13),
\[a_1(\mathbf z)^2A'(\mathbf z)=B'(\mathbf z)b_1(\mathbf z)^2.\]
Since we assumed that $A'$ and $B'$ are independent, this implies that $A'$ and $B'$ are both squares. If $Q'_1\not\equiv0$ (and in particular, $\alpha_1\ne0$), then either $a_1(\mathbf z)=b_1(\mathbf z)\equiv0$, in which case Equation(3.13) implies that $Q'_1(\mathbf z)^2=\alpha_1^2A'(\mathbf z)B'(\mathbf z)$, and we are done (as either both $A'$ and $B'$ are squares or they are both multiples of $Q'_1$), or equations(3.14) and(3.15) imply that $\alpha_1^2A'(\mathbf z)B'(\mathbf z)=Q'_1(\mathbf z)^2$ which again implies the claim.
This concludes the proof of Theorem3.1 for the case $b\equiv0$ and thus the proof of the theorem.
\end{proof}
